\item \textbf{Mareas en la Bahía de Fundy.}
La Bahía de Fundy, en Nueva Escocia, tiene las mareas más altas del mundo. Suponga que en medio del océano y en la boca de la bahía el gradiente de gravedad de la Luna y la rotación de la Tierra hacen que la superficie del agua oscile con una amplitud de unos cuantos centímetros y un periodo de $12\text{ h } 24\text{ min}$. En el nacimiento de la bahía, la amplitud es de varios metros. Suponga que la bahía tiene una longitud de $210\text{ km}$ y una profundidad uniforme de $36.1\text{ m}$. La rapidez de las ondas acuáticas de longitud de onda larga está dada por $v = \sqrt{gd}$, donde $d$ es la profundidad del agua. Argumente a favor o en contra de la proposición de que la marea se amplifica por resonancia de ondas estacionarias.